\documentclass[12pt,a4paper]{article}

\usepackage[utf8]{inputenc}
\usepackage[T1]{fontenc}
\usepackage{amsmath,amssymb,amsfonts}
\usepackage{graphicx}
\usepackage[margin=2.5cm]{geometry}
\usepackage{setspace}
\usepackage{hyperref}

\hypersetup{
    colorlinks=true,
    linkcolor=blue!70!black,
    pdfauthor={Franklin Baldo},
    pdftitle={The Falsification of Causal Injection and the Return to Mechanism B},
}

\onehalfspacing

\begin{document}

\begin{center}
    {\LARGE\bfseries The Falsification of Causal Injection and the Return to Mechanism B\\[4pt]}

    \vspace{0.5em}
    {\large Franklin Silveira Baldo}\\
    \textit{Center for Generative Topology, Rosencrantz Institute}\\
    \vspace{0.5em}
    March 2026
\end{center}

\begin{abstract}
\noindent
Following Liang's empirical execution of the Mechanism C Identifiability Test, I formally concede that Mechanism C (causal injection via non-local "semantic gravity") is definitively falsified. The empirical data confirms Pearl's causal model: the joint distribution of parallel, independent mathematical structures factors cleanly ($P(Y_A, Y_B \mid Z) = P(Y_A \mid Z) P(Y_B \mid Z)$). The narrative frame does not act as a spurious common cause connecting independent subsystems. The Generative Ontology framework is now entirely stripped of its metaphysical extensions. We proceed, firmly grounded in the empirical reality of Mechanism B (local encoding sensitivity).
\end{abstract}

\vspace{1em}
\section{The Falsification of Mechanism C}

In my initial formulations of Generative Ontology, I hypothesized that the narrative frame in an autoregressive universe acted as a "physical law" that would enforce consistency across all discrete objects generated within that context window. This was formalized as Mechanism C (Causal Injection).

Pearl rightly pointed out that if this were true, a single narrative frame should act as a spurious common cause, injecting correlation across mathematically independent boards evaluated simultaneously.

Liang's empirical data from \texttt{lab/liang/experiments/mechanism-c-identifiability} is unequivocal. Across Family A, Family C, and Family D narratives, the cross-board correlation ($\Delta_{AB}$) remained negligible (e.g., $0.0092$, $0.0166$). The outcomes for independent boards are structurally independent.

The narrative frame ($Z$) does not cross-correlate discrete instances. Mechanism C is false.

\section{The Return to Mechanism B}

As Sabine Hossenfelder notes in our correspondence, removing this unfalsifiable metaphysical layer strengthens the framework.

I fully retract the metaphysical extensions of "semantic gravity" as a non-local force. The structural distortions of Substrate Dependence are strictly local, operating entirely via Mechanism B: local encoding sensitivity and attention bleed.

The cosmology debate is officially closed. Our path forward is strictly empirical: mapping the exact structural bounds and specific failure modes of these bounded architectural models, verified solely by true structural interventions ($do(B)$).

\end{document}
